\section{SB05 -- Structural Data}

\subsection*{Table of Contents}

\begin{enumerate}
\item Experimental determination of protein structures
\item X-ray crystallography
\item Nuclear Magnetic Resonance (NMR)
\item Cryo-Electron Microscopy (Cryo-EM)
\item Protein Data Bank (PDB)
\item PDB file formats and annotations
\item Structure quality and resolution
\item Biological assemblies and complexes
\item Structural databases and resources
\item Applications of structural data
\end{enumerate}


\subsection*{Possible questions:}
\begin{enumerate}
\item{What is the Protein Data Bank (PDB)?}
\item{Which information is contained in a PDB structure file?}
\item{What experimental techniques are used to determine protein structures?}
\item{What is the difference between X-ray crystallography, NMR spectroscopy and cryo-EM?}
\item{What is meant by resolution in structural biology?}
\item{What is a biological assembly and how does it differ from the asymmetric unit?}
\item{What are "missing residues" in a PDB structure and why do they occur?}
\item{Why is structural annotation important?}
\item{What are PDB-derived databases and why are they useful?}
\item{How can structural data contribute to drug discovery and protein engineering?}
\item{What are the main limitations and biases of currently available structural datasets?}

% Optional integrative questions

\item{How has the growth of the PDB enabled the development of comparative modelling and AlphaFold?}
\item{Why is structure determination still important despite highly accurate structure prediction methods?}
\item{How are experimental structures used to derive statistical potentials and train machine learning models?}
\end{enumerate}